% FORMALIZACION_NUEVA: explicit proof of decomposition and
% irreducibility already stated in the five-component development of REV06.
% CERTIFICADO_NUEVO: verificar_caracteres_pentadicos_rev07.py.
% Intended insertion after the definition of V12, before its decomposition.
\paragraph{Characters of icosahedral incidence.}
The following calculation is performed on the icosahedral action obtained in
Theorem~\ref{ii:pent:thm:gravity-exceptional-descent}, with its incidence
and orientation hypotheses. It determines the decomposition of the
vertex representation and the irreducibility of the tensor space
through their traces, preserving the antecedent
\(\pi_{\mathrm{ico}}\) declared in that theorem.

\begin{proposition}[Vertex decomposition and tensor character]
\label{prop:ii-pent-caracteres}
Let \(\tau=(1+\sqrt5)/2\) and
\(\bar\tau=(1-\sqrt5)/2\). In the class order
\(1,2A,3A,5A,5B\), the characters of the rotational representation
\(V_3\), its algebraic conjugate \(V_{3'}\), the twelve-vertex
representation \(V_{12}\), and
\(\operatorname{STF}_3=\operatorname{Sym}^2_0(V_3)\) are
\begin{equation}
\label{eq:ii-pent-tabla-caracteres}
\begin{array}{c|rrrrr}
\text{class} &1&2A&3A&5A&5B\\ \hline
\text{cardinality}&1&15&20&12&12\\
\text{class of }g^2&1&1&3A&5B&5A\\ \hline
\chi_3&3&-1&0&\tau&\bar\tau\\
\chi_{3'}&3&-1&0&\bar\tau&\tau\\
\chi_{12}&12&0&0&2&2\\
\chi_{\mathrm{STF}}&5&1&-1&0&0
\end{array}
\end{equation}
The representations \(V_3,V_{3'}\) and
\(\operatorname{STF}_3\) are absolutely irreducible and distinct.
Moreover,
\begin{equation}
\label{eq:ii-pent-descomposicion-demostrada}
 V_{12}\simeq\mathbf1\oplus V_3\oplus V_{3'}
                  \oplus\operatorname{STF}_3,
 \qquad
 \langle\chi_{12},\chi_{\mathrm{STF}}\rangle=1.
\end{equation}
In particular, the five-dimensional summand occurs exactly once.
\end{proposition}

\begin{proof}
In \(A_5\), there are \(5\cdot3=15\) double transpositions and
\(\binom53\cdot2=20\) three-cycles.
The \(4!=24\) five-cycles split into two classes of
twelve: their centralizer is the cyclic subgroup of order five and is
contained in \(A_5\). Identifying the positions of a cycle
with \(\mathbb Z/5\mathbb Z\), the permutation inducing
\(r\mapsto2r\) is a four-cycle on the nonzero
elements, with negative sign. Squaring therefore exchanges the two
classes. Inversion \(r\mapsto-r\) is a product of two
transpositions and preserves each class. This determines the first two
rows of~\eqref{eq:ii-pent-tabla-caracteres}.

The representation \(V_3\) is the rotation action on the
vertex coordinates. A rotation through angle \(\theta\)
has trace \(1+2\cos\theta\). The classes of orders two and three
give \(-1\) and \(0\), respectively. The class of angles \(\pm2\pi/5\)
is called \(5A\), and that of
\(\pm4\pi/5\) is called \(5B\); their traces are \(\tau\) and \(\bar\tau\).
These expressions follow from
\(u^2+u-1=0\), for \(u=2\cos(2\pi/5)\), and the double-angle
formula. The equation for \(u\) follows by dividing
\(1+z+z^2+z^3+z^4=0\) by \(z^2\), with
\(z=\exp(2\pi i/5)\). The character of an already geometrically
defined representation has thus been calculated.

To construct the other representation, use the vertices without the
common normalization \(\nu^{-1}\). Their coordinates belong to
\(K=\mathbb Q(\sqrt5)\). If \(B\) is the matrix of three linearly
independent vertices and \(B_g\) contains their images, the
rotation matrix is \(R_g=B_gB^{-1}\in M_3(K)\).
The conjugation \(\sigma:\sqrt5\mapsto-\sqrt5\) then defines
\(R'_g=\sigma(R_g)\), satisfying
\[
 R'_{gh}=R'_gR'_h,\qquad
 (R'_g)^TR'_g=I,\qquad \det R'_g=1.
\]
Thus \(V_{3'}\) is an effectively constructed real
representation. Its character is \(\sigma(\chi_3)\), exchanging
\(\tau\) and \(\bar\tau\).

The permutation character counts fixed vertices. The identity fixes
twelve; nontrivial rotations of order five fix the two endpoints
of their axis. Those of order two or three fix no vertices, since the
stabilizer of a vertex is \(C_5\). Hence
\(\chi_{12}=(12,0,0,2,2)\).

For the tensor space, if \(\lambda_1,\lambda_2,\lambda_3\)
are the eigenvalues of \(R_g\), the sum
\(\sum_{i\leq j}\lambda_i\lambda_j\) is
\(\bigl((\sum_i\lambda_i)^2+\sum_i\lambda_i^2\bigr)/2\).
The metric form generates an invariant line in
\(\operatorname{Sym}^2(V_3)\), whose trace-free complement has
character
\begin{equation}
\label{eq:ii-pent-caracter-stf}
 \chi_{\mathrm{STF}}(g)
 =\frac{\chi_3(g)^2+\chi_3(g^2)}2-1.
\end{equation}
Using the squaring map in the table and
\(\tau+\bar\tau=1\), \(\tau\bar\tau=-1\), gives
\(\chi_{\mathrm{STF}}=(5,1,-1,0,0)\).

The real-character inner product is
\(\langle\chi,\psi\rangle
=\frac1{60}\sum_C|C|\chi(C)\psi(C)\).
In particular,
\begin{align*}
 \langle\chi_3,\chi_3\rangle
 &=\langle\chi_{3'},\chi_{3'}\rangle
 =\frac{9+15+12(\tau^2+\bar\tau^2)}{60}=1,\\
 \langle\chi_3,\chi_{3'}\rangle
 &=\frac{9+15+24\tau\bar\tau}{60}=0,\\
 \langle\chi_{\mathrm{STF}},\chi_{\mathrm{STF}}\rangle
 &=\frac{25+15+20}{60}=1.
\end{align*}
Norm one proves irreducibility of the complexifications;
consequently, it also proves absolute irreducibility of the constructed
real representations. Their dimensions and orthogonality distinguish the
three summands. Finally, the equality of class functions
\[
 \chi_{12}=1+\chi_3+\chi_{3'}+\chi_{\mathrm{STF}}
\]
is verified in all five columns. Semisimplicity of finite-group
representations in characteristic zero gives
decomposition~\eqref{eq:ii-pent-descomposicion-demostrada}. The
multiplicity of the tensor summand can be checked directly:
\[
 \langle\chi_{12},\chi_{\mathrm{STF}}\rangle
 =\frac{12\cdot5}{60}=1.
\]
\end{proof}

Thus the central idempotent associated with
\(\chi_{\mathrm{STF}}=\chi_5\) has rank
\(5\cdot1=5\) in \(V_{12}\). The irreducibility entering
normalization of the tensor intertwiner and the multiplicity
entering gravitational coupling are established by
this same calculation.
